\documentclass[12pt]{article}

% ------------------------------------------------
% PAQUETES
% ------------------------------------------------
\usepackage[spanish,es-noshorthands]{babel}
\usepackage[utf8]{inputenc}
\usepackage[T1]{fontenc}
\usepackage{lmodern}
\usepackage{amsmath,amssymb,amsthm}
\usepackage{array}
\usepackage{booktabs}
\usepackage{enumitem}
\usepackage{geometry}
\usepackage{setspace}

% ------------------------------------------------
% CONFIGURACIÓN
% ------------------------------------------------
\geometry{margin=1.5cm}
\setstretch{1.15}

% ------------------------------------------------
% DOCUMENTO
% ------------------------------------------------
\begin{document}
	
	\begin{center}
		{\Large \textbf{Solución Problema 2}}\\[0.3cm]
		{\normalsize Pensamiento Matemático}\\[0.2cm]
		{\normalsize Contexto: Psicología}
	\end{center}
	
	\vspace{0.4cm}
	
	\noindent
	\textbf{Problema 2}
	
	\vspace{0.2cm}
	
	Considere la siguiente proposición compuesta:
	
	\begin{quote}
		``Si una persona no sigue las estrategias terapéuticas indicadas, entonces presenta dificultades emocionales o no logra mejorar su bienestar psicológico.''
	\end{quote}
	
	\begin{enumerate}[label=\alph*)]
		
		\item \textbf{Identificación de proposiciones simples y expresión en lenguaje simbólico}
		
		Definimos las siguientes proposiciones simples:
		
		$p:\ \text{La persona sigue las estrategias terapéuticas indicadas.}$


		$q:\ \text{La persona presenta dificultades emocionales.}$
		
		$r:\ \text{La persona logra mejorar su bienestar psicológico.}$
		
		
		Ahora traducimos el enunciado parte por parte:
		
		\begin{itemize}
			\item ``La persona no sigue las estrategias terapéuticas indicadas'' se representa por:
			\[
			\sim p
			\]
			
			\item ``presenta dificultades emocionales'' se representa por:
			\[
			q
			\]
			
			\item ``no logra mejorar su bienestar psicológico'' se representa por:
			\[
			\sim r
			\]
			
			\item La expresión ``dificultades emocionales o no logra mejorar su bienestar psicológico'' se representa por:
			\[
			q \lor \sim r
			\]
			
			\item Como todo el enunciado tiene la forma ``Si \dots entonces \dots'', la proposición completa queda:
			\[
			\sim p \to (q \lor \sim r)
			\]
		\end{itemize}
		
		\textbf{Respuesta final:}
		
		La proposición en lenguaje simbólico es:
		\[
		\boxed{\sim p \to (q \lor \sim r)}
		\]
		
		\vspace{0.4cm}
		
		\item \textbf{Tabla de verdad y clasificación lógica}
		
		Construimos la tabla de verdad de la proposición:
		\[
		\sim p \to (q \lor \sim r)
		\]
		
		\textbf{Paso 1:} Recordamos que:
		\begin{itemize}
			\item La negación cambia \(V\) por \(F\) y \(F\) por \(V\).
			\item La disyunción \(q \lor \sim r\) es verdadera cuando al menos una de las dos proposiciones es verdadera.
			\item La implicación \(A \to B\) solo es falsa cuando \(A\) es verdadera y \(B\) es falsa.
		\end{itemize}
		
		\vspace{0.3cm}
		
		\begin{center}
			\begin{tabular}{|c |c| c| c| c| c|}
				\toprule
				\(p\) & \(q\) & \(r\) & \(\sim p\) & \(q \lor \sim r\) & \(\sim p \to (q \lor \sim r)\) \\
				\midrule
				V & V & V & F & V & V \\
				V & V & F & F & V & V \\
				V & F & V & F & F & V \\
				V & F & F & F & V & V \\
				F & V & V & V & V & V \\
				F & V & F & V & V & V \\
				F & F & V & V & F & F \\
				F & F & F & V & V & V \\
				\bottomrule
			\end{tabular}
		\end{center}
		
		\vspace{0.2cm}
		
		\textbf{Paso 2:} Analizamos la última columna.
		
		Observamos que en la columna final aparecen valores verdaderos y también aparece un valor falso.
		
		Por lo tanto:
		
		\begin{itemize}
			\item No es \textbf{tautología}, porque no todos los valores son verdaderos.
			\item No es \textbf{contradicción}, porque no todos los valores son falsos.
			\item Sí es una \textbf{contingencia}, porque tiene al menos un valor verdadero y al menos un valor falso.
		\end{itemize}
		
		\textbf{Justificación puntual del caso falso:}
		
		La proposición resulta falsa cuando:
	$
		p=F,\qquad q=F,\qquad r=V
	$
		
	
		
		\textbf{Respuesta final:}
		
		La proposición
		\[
		\sim p \to (q \lor \sim r)
		\]
		es una:
		\[
		\boxed{\text{Contingencia}}
		\]
		
	\end{enumerate}
	
\end{document}